% =============================================================================
% 3. Fundamentação Teórica
% =============================================================================

\section{Fundamentação Teórica}

Esta seção apresenta os conceitos fundamentais que sustentam o estudo experimental proposto. A discussão parte do \textit{Vector Space Model} clássico, atravessa a evolução para \textit{embeddings} densos gerados por modelos neurais, descreve algoritmos de busca aproximada de vizinhos próximos em espaços de alta dimensão, caracteriza a arquitetura de bancos de dados vetoriais e contextualiza as metodologias de \textit{benchmarking} aplicáveis.

\subsection{Vector Space Model e Recuperação de Informação}

O \textit{Vector Space Model} (VSM), originalmente proposto por \citeonline{salton1975vsm}, representa documentos e consultas como vetores em um espaço multidimensional, onde cada dimensão corresponde a um termo do vocabulário. A relevância de um documento para uma consulta é estimada pela proximidade angular entre os vetores, tipicamente medida pela similaridade de cosseno. Em sua formulação clássica, os pesos das dimensões são calculados com esquemas como TF-IDF (\textit{term frequency--inverse document frequency}), que combinam a frequência local de um termo com sua raridade no \textit{corpus} \cite{paiva2014semrelext}.

Embora o VSM clássico ofereça uma base sólida para recuperação de informação baseada em palavras explícitas, ele é limitado quanto a capturar relações semânticas implícitas entre termos. \citeonline{costa2014sense} propôs o \textit{framework} SENSE (\textit{Semantic Enrichment kNowledge SourcEs}), que estende o VSM com suporte ontológico, incorporando relações entre conceitos de uma ontologia de domínio à representação vetorial. O \textit{framework} demonstra que a qualidade da recuperação aumenta quando a representação captura mais do que apenas a ocorrência de palavras explícitas. As representações vetoriais densas adotadas hoje em sistemas de busca semântica e RAG podem ser entendidas como a evolução natural dessa linha de raciocínio: substituem a engenharia manual de ontologias por aprendizado automático de relações semânticas a partir de grandes \textit{corpora} textuais.

\todo{Expandir esta subseção com a descrição formal de TF-IDF, similaridade de cosseno e limitações do VSM clássico (sinonímia, polissemia, vocabulary mismatch).}

\subsection{Embeddings Densos e Modelos de Linguagem}

Embora o VSM clássico ofereça uma representação vetorial bem definida, sua dimensionalidade é determinada pelo tamanho do vocabulário e seus vetores são esparsos: cada documento ativa um pequeno subconjunto das dimensões do espaço, e termos sinônimos ou semanticamente relacionados ocupam dimensões independentes. Essa propriedade dificulta a captura de relações de significado que não estão expressas pela coincidência exata de palavras.

Modelos neurais de linguagem permitiram a transição para representações vetoriais densas, em que cada documento ou sentença é mapeado para um espaço de algumas centenas a poucos milhares de dimensões, todas ativas. Nessa representação, vetores de textos com significados próximos tendem a ocupar regiões próximas do espaço, mesmo na ausência de palavras em comum. \citeonline{reimers2019sbert} propuseram a arquitetura \textit{Sentence-BERT}, que adapta modelos do tipo BERT por meio de redes siamesas para gerar \textit{embeddings} de sentenças cuja similaridade pode ser comparada diretamente por produto interno ou similaridade de cosseno, viabilizando o uso desses vetores em tarefas de recuperação semântica em larga escala.

A dimensionalidade dos \textit{embeddings} produzidos por modelos contemporâneos varia tipicamente entre 384 e 1536 dimensões, sendo um parâmetro de impacto direto neste estudo: maior dimensionalidade tende a aumentar a expressividade da representação, mas eleva o custo de armazenamento, o \textit{footprint} de memória do índice e o custo computacional das operações de similaridade. A escolha do modelo de \textit{embedding} é, portanto, uma variável metodológica relevante e foi fixada nos experimentos para isolar o efeito do banco de dados sobre o desempenho.

\subsection{Busca Aproximada de Vizinhos Próximos (ANN)}

Dada uma coleção de N vetores em um espaço de alta dimensionalidade, a busca exata pelos K vizinhos mais próximos de um vetor de consulta exige, em geral, comparação contra todos os N vetores, com custo linear na coleção. Esse custo torna-se proibitivo em aplicações de busca semântica que exigem latência interativa sobre coleções com centenas de milhares ou milhões de vetores. Para contornar essa limitação, sistemas modernos adotam algoritmos de busca aproximada de vizinhos próximos (\textit{Approximate Nearest Neighbor}, ANN), que renunciam à garantia de retornar exatamente os K vizinhos mais próximos em troca de ganhos expressivos de tempo de resposta, mantendo \textit{recall} elevado em relação à busca exata.

O algoritmo \textit{Hierarchical Navigable Small World} (HNSW), proposto por \citeonline{malkov2018hnsw}, tornou-se o índice predominante em bancos de dados vetoriais por oferecer um equilíbrio favorável entre \textit{recall}, latência e consumo de memória. O HNSW organiza os vetores em uma estrutura de grafos em múltiplos níveis: camadas superiores contêm poucos nós e arestas longas, permitindo que a busca avance rapidamente entre regiões distantes do espaço; camadas inferiores são mais densas e refinam a busca em torno do candidato. Os parâmetros principais que controlam o \textit{trade-off} entre qualidade e custo são \textit{M} (número de conexões mantidas por nó) e \textit{ef} (tamanho da fila de candidatos durante a busca).

Os três sistemas avaliados neste estudo --- pgvector, Qdrant e Weaviate --- oferecem o HNSW como índice principal. Isso significa que diferenças observadas entre eles refletem decisões de implementação, integração com o resto do SGBD, gerenciamento de recursos e suporte a operações auxiliares (filtros, atualizações, persistência), e não diferenças fundamentais no algoritmo de busca em si. Essa observação é relevante para a interpretação dos resultados: o estudo compara sistemas que partilham o mesmo núcleo algorítmico, mas o aplicam em arquiteturas distintas.

\subsection{Bancos de Dados Vetoriais}

\citeonline{pan2023vdbms} caracterizam os \textit{Vector Database Management Systems} (VDBMS) como sistemas projetados para armazenar, indexar e consultar coleções de vetores densos com suporte a operações de busca por similaridade em escala. Os autores agrupam as soluções existentes em duas categorias arquiteturais principais: (i) sistemas especializados, projetados desde o início para vetores densos e busca aproximada, com índices nativos e otimizações específicas para esse tipo de carga; e (ii) extensões de SGBDs existentes, que adicionam tipos de dados vetoriais e operadores de similaridade a bancos relacionais ou de documentos preexistentes.

O PostgreSQL com a extensão pgvector representa a segunda categoria: estende um SGBD relacional maduro com um tipo de dado vetorial e índices ANN (incluindo HNSW e IVFFlat), permitindo que vetores e metadados estruturados convivam no mesmo banco. Essa abordagem favorece a simplicidade operacional, especialmente em projetos que já operam sobre PostgreSQL, e permite que filtros relacionais sejam combinados com busca vetorial dentro do mesmo plano de execução.

Qdrant e Weaviate representam a primeira categoria. Qdrant é um banco vetorial \textit{open source} escrito em Rust, com suporte nativo a HNSW, filtros por \textit{payload} arbitrário e API REST/gRPC. Weaviate, também \textit{open source}, oferece uma camada de esquema com tipagem, suporte a GraphQL e integração com modelos de \textit{embedding}. Ambos foram projetados desde a origem em torno do problema de recuperação vetorial e tendem a oferecer maior controle sobre parâmetros de indexação e estratégias de filtragem específicas para esse contexto.

A combinação de busca vetorial com filtros estruturados --- central ao Cenário B deste estudo --- é um eixo no qual os três sistemas adotam estratégias distintas, tipicamente classificadas como pré-filtragem, pós-filtragem e filtragem \textit{inline} durante a navegação no grafo HNSW. \citeonline{patel2024acorn} sistematizam o \textit{trade-off} entre essas estratégias e demonstram, com a proposta ACORN, ganhos expressivos de \textit{throughput} sob \textit{recall} fixo quando a filtragem é integrada ao percurso do índice. Essa caracterização fornece base para interpretar diferenças observadas entre pgvector, Qdrant e Weaviate quando consultas vetoriais são combinadas com predicados sobre metadados.

A decisão entre uma solução integrada e uma especializada envolve um \textit{trade-off} explícito: a primeira reduz a complexidade de \textit{stack} ao custo de potencial perda de desempenho em cargas vetoriais intensivas; a segunda otimiza para o caso vetorial mas exige operar um sistema adicional, com seu próprio modelo de dados, ferramental e perfil operacional. Caracterizar esse \textit{trade-off} de forma quantitativa, em cenários representativos, é o objetivo central deste estudo.

\subsection{Retrieval-Augmented Generation (RAG)}

Sistemas \textit{Retrieval-Augmented Generation}, propostos por \citeonline{lewis2020rag}, combinam um modelo de linguagem com um mecanismo de recuperação de informação para gerar respostas fundamentadas em uma base documental externa. O fluxo típico converte a consulta do usuário em um \textit{embedding}, recupera do banco vetorial os trechos de documento mais similares e injeta esse contexto na chamada ao LLM, que gera a resposta final. Essa arquitetura permite que o modelo acesse conhecimento atualizado ou específico de domínio sem necessidade de retreinamento, reduzindo alucinações e aumentando a precisão factual \cite{jing2024llmvdb}.

A qualidade da resposta gerada por um sistema RAG depende criticamente da qualidade da etapa de recuperação: contextos irrelevantes ou incompletos comprometem o resultado mesmo quando o gerador é capaz. \citeonline{pawlik2025rag} demonstra que a configuração e o \textit{tuning} da camada vetorial impactam diretamente a precisão e a robustez de sistemas RAG, e \citeonline{bovas2025navi}, em uma implementação aplicada a contexto educacional, observam \textit{scores} de relevância na faixa de 0,7 a 0,85 condicionados à adequação dessa camada. Esses resultados motivam diretamente o foco deste estudo: avaliar de forma sistemática o desempenho da camada vetorial sob cargas representativas de aplicações RAG.

\subsection{Metodologias de Benchmarking}

O \textit{benchmarking} sistemático de sistemas de banco de dados é uma prática consolidada na literatura de SGBD, e a sua aplicação a sistemas vetoriais demanda atenção a dois eixos complementares: o desempenho operacional clássico (latência, \textit{throughput}, \textit{footprint}) e a qualidade da recuperação (\textit{recall} em relação à busca exata). \citeonline{aerospike2025} sistematiza boas práticas gerais para \textit{benchmarking} de bancos de dados, com ênfase na definição clara de \textit{workloads}, no controle do ambiente e na medição reproduzível de métricas --- princípios diretamente aplicáveis a este estudo.

Para sistemas baseados em ANN, \citeonline{aumuller2020annbench} propuseram o ANN-Benchmarks, uma metodologia que tornou-se referência para avaliação comparativa de algoritmos e implementações de busca aproximada. O ANN-Benchmarks fixa \textit{datasets} de \textit{embeddings}, varre o espaço de parâmetros de cada algoritmo e reporta a curva \textit{recall} $\times$ QPS, evidenciando o \textit{trade-off} fundamental entre qualidade da recuperação e taxa de consultas. Esta metodologia inspira diretamente as decisões deste trabalho: cada sistema é avaliado em um conjunto de configurações de parâmetros do índice HNSW, e os resultados são reportados em curvas que explicitam o \textit{trade-off recall} $\times$ latência, em vez de números pontuais.

As métricas adotadas combinam, portanto, instrumentação clássica de SGBD (latência p50/p95/p99, \textit{throughput} sob carga controlada, uso de memória e disco, tempo de indexação) com métricas específicas de qualidade de recuperação vetorial (recall@K em relação a uma busca exata como referência), aplicadas em três cenários que reproduzem o perfil de uso de aplicações de busca semântica e RAG.
